% GENERATED FILE. Edit canonical lessons, then run npm run generate:books.
% canonical-source-hash: fnv1a64:d8f5d41c30b0d3a3
% canonical-lessons: JA-C86-saku, JA-C86-toku, JA-C86-daku, JA-C86-fumu, JA-C86-kamu

\chapter{To bloom, To solve, To hold in one's arms, To step on, To bite}
\label{ch:ja-to-bloom-to-solve-to-hold-in-one-s-arms-to-step-on-to-bite}
\hlchaptermodality{86}

\begin{chapteropening}
\textbf{By the end of this chapter:} \emph{I can say to bloom, to solve, to hold in one's arms, to step on and to bite.}

Complete the last lesson of chapter 86: i can say to bloom, to solve, to hold in one's arms, to step on and to bite.
\end{chapteropening}

\section[\texorpdfstring{\ja{さく} (saku)}{さく (saku)}]{\ja{さく} (saku) — to bloom}

\label{lesson:JA-C86-saku}



\begin{quote}
Before the new one: say the Japanese for to cook, then the Japanese for to put on.
\end{quote}

\subsection*{You'll want to know: \ja{さく}}
\textbf{\ja{さく}} — \emph{saku} — \textquotedblleft{}to bloom\textquotedblright{}.

\begin{grammarlens}[title={What you've built}]
One more word for this chapter.
\end{grammarlens}

\subsection*{Your turn}
\begin{itemize}
\raggedright
  \item \emph{Say it:} \emph{saku}
  \item \emph{Say it:} \emph{saku}, once more
  \item \emph{Say it:} say \emph{saku}
  \item \emph{Recall:} say the Japanese for to cook, then the Japanese for to put on, then say \emph{saku} again
\end{itemize}

\begin{tcolorbox}[breakable,colback=teal!4,colframe=teal!35!black,title={Before you move on}]
What does \ja{さく} mean? (\textquotedblleft{}To bloom\textquotedblright{}.) Say it once more.
\end{tcolorbox}
\section[\texorpdfstring{\ja{とく} (toku)}{とく (toku)}]{\ja{とく} (toku) — to solve}

\label{lesson:JA-C86-toku}



\begin{quote}
Before the new one: say the Japanese for to put on, then the Japanese for to bloom.
\end{quote}

\subsection*{You'll want to know: \ja{とく}}
\textbf{\ja{とく}} — \emph{toku} — \textquotedblleft{}to solve\textquotedblright{}.

\begin{grammarlens}[title={What you've built}]
One more word for this chapter.
\end{grammarlens}

\subsection*{Your turn}
\begin{itemize}
\raggedright
  \item \emph{Say it:} \emph{toku}
  \item \emph{Say it:} \emph{toku}, once more
  \item \emph{Say it:} say \emph{toku}
  \item \emph{Recall:} say the Japanese for to put on, then the Japanese for to bloom, then say \emph{toku} again
\end{itemize}

\begin{tcolorbox}[breakable,colback=teal!4,colframe=teal!35!black,title={Before you move on}]
What does \ja{とく} mean? (\textquotedblleft{}To solve\textquotedblright{}.) Say it once more.
\end{tcolorbox}
\section[\texorpdfstring{\ja{だく} (daku)}{だく (daku)}]{\ja{だく} (daku) — to hold in one's arms}

\label{lesson:JA-C86-daku}



\begin{quote}
Before the new one: say the Japanese for to bloom, then the Japanese for to solve.
\end{quote}

\subsection*{You'll want to know: \ja{だく}}
\textbf{\ja{だく}} — \emph{daku} — \textquotedblleft{}to hold in one's arms\textquotedblright{}.

\begin{grammarlens}[title={What you've built}]
One more word for this chapter.
\end{grammarlens}

\subsection*{Your turn}
\begin{itemize}
\raggedright
  \item \emph{Say it:} \emph{daku}
  \item \emph{Say it:} \emph{daku}, once more
  \item \emph{Say it:} say \emph{daku}
  \item \emph{Recall:} say the Japanese for to bloom, then the Japanese for to solve, then say \emph{daku} again
\end{itemize}

\begin{tcolorbox}[breakable,colback=teal!4,colframe=teal!35!black,title={Before you move on}]
What does \ja{だく} mean? (\textquotedblleft{}To hold in one's arms\textquotedblright{}.) Say it once more.
\end{tcolorbox}
\section[\texorpdfstring{\ja{ふむ} (fumu)}{ふむ (fumu)}]{\ja{ふむ} (fumu) — to step on}

\label{lesson:JA-C86-fumu}



\begin{quote}
Before the new one: say the Japanese for to solve, then the Japanese for to hold in one's arms.
\end{quote}

\subsection*{You'll want to know: \ja{ふむ}}
\textbf{\ja{ふむ}} — \emph{fumu} — \textquotedblleft{}to step on\textquotedblright{}.

\begin{grammarlens}[title={What you've built}]
One more word for this chapter.
\end{grammarlens}

\subsection*{Your turn}
\begin{itemize}
\raggedright
  \item \emph{Say it:} \emph{fumu}
  \item \emph{Say it:} \emph{fumu}, once more
  \item \emph{Say it:} say \emph{fumu}
  \item \emph{Recall:} say the Japanese for to solve, then the Japanese for to hold in one's arms, then say \emph{fumu} again
\end{itemize}

\begin{tcolorbox}[breakable,colback=teal!4,colframe=teal!35!black,title={Before you move on}]
What does \ja{ふむ} mean? (\textquotedblleft{}To step on\textquotedblright{}.) Say it once more.
\end{tcolorbox}
\section[\texorpdfstring{\ja{かむ} (kamu)}{かむ (kamu)}]{\ja{かむ} (kamu) — to bite, to chew}

\label{lesson:JA-C86-kamu}



\begin{quote}
Before the new one: say the Japanese for to hold in one's arms, then the Japanese for to step on.
\end{quote}

\subsection*{You'll want to know: \ja{かむ}}
\textbf{\ja{かむ}} — \emph{kamu} — \textquotedblleft{}to bite, to chew\textquotedblright{}.

\begin{grammarlens}[title={What you've built}]
One more word for this chapter.
\end{grammarlens}

\subsection*{Your turn}
\begin{itemize}
\raggedright
  \item \emph{Say it:} \emph{kamu}
  \item \emph{Say it:} \emph{kamu}, once more
  \item \emph{Say it:} say \emph{kamu}
  \item \emph{Recall:} say the Japanese for to hold in one's arms, then the Japanese for to step on, then say \emph{kamu} again
\end{itemize}

\begin{tcolorbox}[breakable,colback=teal!4,colframe=teal!35!black,title={Before you move on}]
What does \ja{かむ} mean? (\textquotedblleft{}To bite, to chew\textquotedblright{}.) Say it once more.
\end{tcolorbox}
